% ============================================================
% ENTORNOS PARA CONTRIBUCIÓN POR ESTUDIANTE
% ============================================================

\usepackage{tcolorbox}
\tcbuselibrary{breakable,skins}

% ------------------------------------------------------------
% COLORES PARA TAREAS COMPLETADAS
% ------------------------------------------------------------
\definecolor{donebg}{RGB}{232, 245, 233}         % Fondo verde muy claro
\definecolor{doneborder}{RGB}{21, 71, 52}        % Verde oscuro
\definecolor{donetitlebg}{RGB}{21, 71, 52}       % Verde oscuro
\definecolor{donetitlefg}{RGB}{255,255,255}      % Blanco

% ------------------------------------------------------------
% COLORES PARA TAREAS EN PROCESO
% ------------------------------------------------------------
\definecolor{progbg}{RGB}{255,235,238}           % Fondo Rojo muy claro
\definecolor{progborder}{RGB}{139,0,0}           % Rojo
\definecolor{progtitlebg}{RGB}{139,0,0}          % Rojo
\definecolor{progtitlefg}{RGB}{255,255,255}      % Blanco

% ============================================================
% CAJA PARA TAREAS COMPLETADAS
% ============================================================
\newtcolorbox{taskdonebox}[2][]{
    enhanced,
    breakable,
    colback=donebg,
    colframe=doneborder,
    coltitle=donetitlefg,
    colbacktitle=donetitlebg,
    boxrule=0.8pt,
    arc=3mm,
    left=4mm,
    right=4mm,
    top=2mm,
    bottom=2mm,
    title={\textbf{#2}},
    fonttitle=\bfseries\sffamily,
    attach boxed title to top left={yshift*=-2mm,xshift=2mm},
    boxed title style={
            size=small,
            sharp corners=north,
            bottom=0.5mm,
            left=2mm,
            right=2mm,
        },
    #1
}

% Entorno más cómodo para tareas completadas
\newenvironment{taskdone}[2][]{
    % #1: opciones extra, #2: título de la tarea
    \begin{taskdonebox}[#1]{#2}
        }{
    \end{taskdonebox}
}

% ============================================================
% CAJA PARA TAREAS EN PROCESO
% ============================================================
\newtcolorbox{taskprogressbox}[2][]{
    enhanced,
    breakable,
    colback=progbg,
    colframe=progborder,
    coltitle=progtitlefg,
    colbacktitle=progtitlebg,
    boxrule=0.8pt,
    arc=3mm,
    left=4mm,
    right=4mm,
    top=2mm,
    bottom=2mm,
    title={\textbf{#2}},
    fonttitle=\bfseries\sffamily,
    attach boxed title to top left={yshift*=-2mm,xshift=2mm},
    boxed title style={
            size=small,
            sharp corners=north,
            bottom=0.5mm,
            left=2mm,
            right=2mm,
        },
    #1
}

% Entorno más cómodo para tareas en proceso
\newenvironment{taskinprogress}[2][]{
    % #1: opciones extra, #2: título de la tarea
    \begin{taskprogressbox}[#1]{#2}
        }{
    \end{taskprogressbox}
}

% ============================================================
% CALLOUTS (Nota, Advertencia, Error)
% ============================================================

% Nota (callout azul)
\newcommand{\tasknote}[1]{%
    \par\medskip
    \noindent\textbf{\color{blue!70!black}Nota:}~#1\par
}

% Advertencia (callout naranja)
\newcommand{\taskwarn}[1]{%
    \par\medskip
    \noindent\textbf{\color{orange!70!black}Advertencia:}~#1\par
}

% Error o problema encontrado (callout rojo)
\newcommand{\taskerror}[1]{%
    \par\medskip
    \noindent\textbf{\color{red!70!black}Error/Problema:}~#1\par
}

% Pendiente (callout morado - útil para taskinprogress)
\newcommand{\taskpending}[1]{%
    \par\medskip
    \noindent\textbf{\color{purple!70!black}Pendiente:}~#1\par
}